% Generated draft: requires mathematical and source review.
\documentclass[11pt]{article}
\usepackage[margin=25mm]{geometry}
\usepackage{fontspec}
\setmainfont{Latin Modern Roman}
\usepackage{amsmath,amssymb,mathtools,unicode-math}
\setmathfont{Latin Modern Math}
\usepackage{enumitem,needspace,xurl,hyperref}
\setlength{\parindent}{0pt}
\setlength{\parskip}{5pt}
\setlength{\emergencystretch}{3em}
\newcommand{\N}{\mathbb{N}}
\newcommand{\Z}{\mathbb{Z}}
\newcommand{\Q}{\mathbb{Q}}
\newcommand{\R}{\mathbb{R}}
\newcommand{\C}{\mathbb{C}}
\newcommand{\F}{\mathbb{F}}
\newcommand{\A}{\mathbb{A}}
\newcommand{\PP}{\mathbb{P}}
\newcommand{\E}{\mathbb{E}}
\newcommand{\eps}{\varepsilon}
\newcommand{\dd}{\,\mathrm{d}}
\begin{document}
\section*{UnsolvedMath ID: 30002177}
\textbf{Source problem code:} catalogue entry
\subsection{Definitions and mathematical statement}
Let \(n\) be a positive integer and let \(f\in\Z[t_1,\ldots,t_n]\). A root of unity is a complex number \(z\) such that \(z^m=1\) for some positive integer \(m\); its order \(\operatorname{ord}(z)\) is the least such integer. Write \(|\cdot|\) for the usual complex modulus.

The supplied question asks whether there exists a real number \(E=E(f)>0\), independent of the order bound and the chosen roots, such that
\[
\forall d\in\Z_{\ge 1}\quad
\forall(z_1,\ldots,z_n)\in\C^n,\qquad
\left(\operatorname{ord}(z_i)\le d\text{ for every }i\right)
\Longrightarrow
\left(f(z_1,\ldots,z_n)=0\ \text{or}\ |f(z_1,\ldots,z_n)|>d^{-E}\right),
\]
where the tuple is restricted to roots of unity.

\emph{Unresolved qualification.} The displayed quantification over positive integers \(d\) makes explicit the natural literal reading of the supplied order bound; the source record does not specify its range. At \(d=1\), all \(z_i=1\), and the strict inequality requires \(|f(1,\ldots,1)|>1\) whenever this value is nonzero. Thus nonzero values of modulus \(1\) are excluded by the literal formulation. The record does not resolve whether the intended statement restricts \(d\), changes the strict inequality, or includes another qualification. No such repair is adopted here.
\subsection{Short English statement}
The supplied conjecture asks for an exponent depending only on the integer polynomial that bounds every nonzero value at roots of unity away from zero by a power of their order bound. The strict inequality and the unspecified range of \(d\), especially the boundary \(d=1\), require clarification.
\subsection{Sources}
\begin{itemize}
\item Local UnsolvedMath registry, permanent ID \(30002177\), title \emph{Polynomial Separation at Roots of Unity}, exact source code \texttt{catalogue entry}. The record supplies proposed year \(2012\) and names no proposer.
\item Supplied report citation: \emph{Low-Dimensional Topology and Number Theory (2013)}, Oberwolfach Reports, pp.\ 2541--2596, DOI \texttt{10.4171/owr/2012/42}; identified in the literature sources as Oberwolfach Report 42, Conjecture 6. \url{https://doi.org/10.4171/owr/2012/42}.
\item Supplied literature reference: Barber, \emph{small sums of five roots of unity}. \url{https://doi.org/10.1112/blms.12826}.
\item Supplied literature reference: Kalogirou, \emph{small sums of roots of unity}. \url{https://arxiv.org/abs/2607.06098}.
\item Literature assessment dated exactly \texttt{2026-08-21}: the supplied curation labels the problem open, states that reduction to sums of roots of unity is known, and reports that no general polynomial separation bound has been proved. This is supplied dated curation, not independently verified. The background explicitly calls for independent verification of time-sensitive status and cited claims.
\item The registry labels the statement \texttt{corrected\_verified} and reports manual semantic, typographic, and LaTeX review. These are supplied editorial assertions, not independent mathematical verification of this draft.
\end{itemize}
\end{document}
